\scanpage{45}
\mspage{\textcircled{1}}
(20.1.3) \inserted{A la place de ``Nous allons nous intéresser... n'est autre que l'ensemble des éléments réguliers de $\mathcal O_{X,x}$'' lire :} ``Nous allons nous intéresser ici au cas où $\mathscr S$ est le sous-faisceau $\mathscr S(\mathcal O_X)$ de $\mathcal O_X$ tel que pour tout ouvert $U$, $\Gamma(U,\mathscr S)$ soit l'ensemble des éléments réguliers de l'anneau $\Gamma(U,\mathcal O_X)$ dont la restriction à tout ouvert $V\subset U$ est encore régulière dans $\Gamma(V,\mathcal O_X)$. Il est immédiat qu'il s'agit d'un faisceau car un élément est dans $\Gamma(U,\mathscr S)$ si et seulement si il est dans $\Gamma(U,\mathcal O_X)$ et \deleted{tous} ses germes en tout point $x$ de $U$ sont des éléments réguliers de $\mathcal O_{X,x}$. Il faut remarquer que $\mathscr S(\mathcal O_X)_x$ n'est pas \deleted{toujo} nécessairement l'ensemble des éléments réguliers de $\mathcal O_{X,x}$.''

(20.1.4) au lieu de : ``qui n'est autre que \deleted{un}\inserted{l'}anneau total des fractions de $\Gamma(U,\mathcal O_X)$'', -- lire : ``qui est un anneau de fractions à dénominateurs dans une partie multiplicativement stable composée d'éléments réguliers''.

(20.1.8) à la place de : ``en effet, si $s\in\Gamma(U,\mathscr S(\unsure{\mathcal M|\mathcal K}_X))$... inversible dans cet anneau de fractions'', -- lire : ``en effet, si $s\in\Gamma(U,\mathscr S(\mathcal M_X))$, pour tout $x\in U$ il existe un voisinage ouvert $V\subset U$ de $x$ tel que $s|V$ soit un élément régulier de $\Gamma(\deleted{\unsure{U}}V,\mathcal O_X)[\Gamma(V,\mathscr S)^{-1}]$ dont toutes les restrictions à des ouverts $W$ inclus dans $V$ sont des éléments réguliers de $\Gamma(W,\mathcal O_X)[\Gamma(W,\mathscr S)^{-1}]$ et on sait que dans ce cas $s|V$ est inversible dans $\Gamma(V,\mathcal O_X)[\Gamma(V,\mathscr S)^{-1}]$.''

(20.1.11) Remplacer ``$\mathscr S_f(U)$ l'ensemble des sections régulières $s\in\Gamma(U,\mathcal O_X)$ telles \inserted{que}...'', -- par : ``$\mathscr S_f(U)$ l'ensemble des sections régulières $s\in\Gamma(U,\mathscr S)$ \deleted{telles que} de $\mathcal O_X$ au dessus de $U$ telles que...''

% archive sheet number: 44

